% !TeX spellcheck = es_ES
\documentclass[11pt,a4paper]{report} % Para tener un documento grande pero sin ser un libro (solo hay esta posibilidad)
\usepackage{../portada}

\title{Adaptaci�n y extensi�n de la herramienta de captura y an�lisis de tr�fico Ksensor}
\author{Javier Domingo Cansino}
\subtitle{Pliego de Condiciones}
\director{Armando Ferro Vazquez}
\proffesor{Luis Zabala Alberdi}
\codirectortrue

\usepackage[latin1]{inputenc}   % Para la codificaci�n
\usepackage{textcomp}           % Para los simbolos de euro
\usepackage[spanish]{babel}     % Para el lenguaje que tenga bien
\usepackage[T1]{fontenc}        %      los saltos de linea con gui�n etc.
\bibliographystyle{alpha}       % Estilo de bibliograf�a
\usepackage{graphicx}           % Para insertar im�genes
\graphicspath{{../images/}}     % La ruta relativa de las im�genes
\usepackage{hyperref}           % Hacer links guays entre secciones y hacia fuera
\usepackage[hypcap]{caption}    % Hacer que las referencias se vean, y no se queden por encima (referenciando imagenes, etc.)
\hypersetup{
	colorlinks=true,			% Los links deben estar coloreados
	linkcolor=blue,				% los links esternos deben ser azules
	citecolor=red,				% y los de las citas rojos
}
\usepackage[table]{xcolor}      % Dar colores a las tablas
\usepackage{etoolbox}			% Para el espaciado
\usepackage[doublespacing]{setspace}
\usepackage{listings}			% Para formatear c�digo
\lstset{						% El estilo para formatearlo sera este
	basicstyle=\small\singlespacing\ttfamily, % La letra tiene que ser peque�a y espaciado simple y monoespaciada
	keywordstyle=\bfseries\color{blue},		% las palabras clave tienen que estar en azul
	identifierstyle=\color{black},	% Los identificadores (funciones y variables)
	commentstyle=\color{gray},		% Los comentarios en gris
	language=C,						% El lenguaje de programaci�n es C
	extendedchars=true,				% Usamos m�s caracteres que los ascii
	breaklines=true,				% Puede cortar las lineas que ocupen demasiado
	breakatwhitespace=true,			% 	pero tiene que hacerlo en los espacios en blanco
	frame=lines,					% Poner lineas al principio y al final del c�digo para acotarlo
	tabsize=4,
}
\renewcommand{\lstlistlistingname}{\textbf{Trozos de c�digo}}
\renewcommand{\lstlistingname}{\textbf{Trozo de c�digo}}
\usepackage{eurosym}            % para tener el simbolo del euro
\usepackage{setspace}           % Para hacer el espaciado entre l�neas
\setstretch{1.4}                %      1.4 del espacio normal
\usepackage{parskip}            % El espaciado entre parrafos 
\setlength{\parindent}{25pt}    %      lo ponemos a 25pt
\usepackage[body=17cm,footskip=2cm,margin=3cm]{geometry} % Cambiar el margen del texto

\usepackage{fancyhdr}           % Cambiar las cabeceras del documento
\pagestyle{fancy}               %      y usar los cambios
\lhead{\includegraphics[width=5cm]{ingenieros-bilbao}}
\rhead{Adaptaci�n y extensi�n de la herramienta\\ de captura y an�lisis de tr�fico Ksensor}
\cfoot{\thepage}
\setlength{\headheight}{46pt}	% Para evitar que desaparezca una imagen

\usepackage{titlesec}           % cambiar el estilo soso de los cap�tulos
\titleformat{\chapter}[hang]{\Huge\bfseries}{\thechapter. }{0pt}{\Huge\bfseries} % Cual es el formato que deben seguir los cap�tulos
\titlespacing{\chapter}{0pt}{-80pt}{20pt} %cual es el espacio que deben tener encima
\usepackage{makeidx}            % Hacer �ndices

% % % % % % % % % % % % % % % % % % % % % % %
% Secci�n de comandos personalizados        %
% % % % % % % % % % % % % % % % % % % % % % %
\newcommand{\reference}[2]{\hyperref[#1]{\textit{\ref*{#1} #2}}}
\newcommand{\code}[1]{\texttt{#1}}
\makeindex

\begin{document}
\pagenumbering{Roman}
\maketitle
\tableofcontents % Crear la tabla de contenidos
\newpage
\pagenumbering{arabic}
% !TeX spellcheck = es_ES
% !TeX root = main.tex
En los últimos años, las tecnologías informáticas de la comunicación están revolucionando la manera de comunicarse del mundo entero, tener acceso a internet se ha convertido en algo necesario para poder ponerse en contacto con el resto del mundo.

Cada vez está todo más interconectado, con el consiguiente crecimiento del tráfico en todas las redes. Incluso empresas que antes contrataban redes privadas para sus comunicaciones, ahora utilizan internet.

Como en cualquier servicio crítico el mantenimiento debe ser proactivo y por ello se han creado métodos para garantizar la máxima eficiencia de las comunicaciones. Para ello se debe garantizar la seguridad en la red y la detección instantánea de problemas, asegurando así la \textit{QoS} (\textit{Quality of Service}, Calidad de Servicio).

Para la securización de la red se implantan \textit{firewalls}, \textit{proxys}, y \textit{IDS}s (\textit{Intrusion Detection System}, Sistema de Detección de Intrusión), que se encargan de evitar accesos externos no autorizados a la red, controlar el tráfico saliente, y detectar posibles atacantes dentro de la red.

El método para asegurar la calidad de las comunicaciones es la medición de parámetros de \textit{QoS}, tales como retardo de los paquetes, errores en la transmisión, tipos de tráfico comunes, rutas óptimas, etc.

Son varias las soluciones comerciales que implementan estos servicios pero actualmente no hay ninguna que implemente todo en un mismo sistema. Además, no hay ningún producto de código libre que aproveche al máximo la posible eficiencia del equipo, haciendo un análisis \textit{online}\footnote{Online se refiere a hacer el análisis mientras se captura, en vez de capturar, guardar y luego analizar.}.

En el grupo de investigación \textit{Network Quality and Security (NQAS)} se ha implementado un sensor que implementa algunas de las citadas funcionalidades, en el diseño se ha contemplado la posibilidad de más adelante poder ampliarlo de una manera sencilla, y como valor añadido fundamental, aprovecha al máximo los recursos disponibles, ya que ha sido programado de una manera que permita ejecución multi-hilo.

Este prototipo se llama ksensor, aunque no es el primero que se ha creado. La línea de investigación principal, llamada hi-sensor, ha ido haciendo cada vez más complejo el diseño y la programación del prototipo con el objetivo de mejorar la eficacia.

A rasgos generales, primero se diseño de forma que la aplicación fuera un programa de usuario que capturaba a través de las llamadas estándar al sistema. Se utilizaron también las librerías para tener computación multi-hilo, y se consiguieron unos resultados prometedores.

Cuando se empezaron a analizar redes de alta velocidad en saturación, y se vio que había un problema: El equipo estaba continuamente capturando, y no había espacio para el análisis. Para solucionar este problema, se acordó migrar a espacio de kernel.

El kernel es la base que hace que todos los programas funcionen. Es un programa que se encarga de gestionar los dispositivos y recursos del ordenador para ofrecer una interfaz libre de detalles a los programas. Además, también se encarga de hacer que varios programas se ejecuten como si se ejecutaran a la vez, que los programas sean capaces de direccionar una memoria que no existe de una forma transparente a ellos, manejando eficazmente una MMU (\textit{Memory Management Unit}, Unidad de Gestión de Memoria), y muchas otras cosas ocultas al programador más.

En general, del diseño de un buen kernel dependerá la eficiencia del sistema. En nuestro caso, esta característica tiene valor añadido porque el diseño del kernel no está hecho para nuestro tipo de sistema. La eficiencia del kernel es buena para los sistemas operativos de carácter general, en los que capturar los paquetes de uno mismo es suficiente, y que es una cantidad pequeña en comparación con el tráfico de la red.

El sistema que tenemos es únicamente para capturar y procesar y por lo tanto, la máxima eficiencia se puede describir como el equilibrio entre el tiempo en el que el ordenador esta capturando y el que está analizando. Se tienen que capturar exactamente el número de paquetes que se van a analizar.

Ahora, el sistema es una
\chapter{CONDICIONES T�CNICAS}

En este apartado se exponen los requerimientos que tanto los recursos materiales como humanos deben cumplir para la correcta realizaci�n del presente proyecto.

\section{Recursos materiales}
El presente proyecto no requiere de recursos materiales hardware espec�ficos, si bien es necesario un ordenador con un procesador multicore a f�n de hacer los tiempos de compilaci�n en desarrollo m�s cortos. En cuanto a los recursos materiales software, se recurre al software libre para cubrir dichas necesidades. As� se utiliza el siguiente software:

Sistema operativo Debian GNU/Linux Wheezy, con los siguiente paquetes instalados:
\begin{itemize}
	\item \textbf{build-essential} Un conjunto de aplicaciones que incluye todo lo b�sico para compilar un programa
	\item \textbf{git} El sistema de control de versiones que se va a utilizar
\end{itemize}

\subsection{Recursos humanos}

Las personas que se encargan del dise�o y desarrollo del presente proyecto deben cumplir con los perfiles aqu� indicados:
\begin{description}
\item[Director del proyecto] Ingeniero de Telecomunicaciones con experiencia en direcci�n de proyectos similares. Ser� el encargado de supervisar los dise�os y de que se cumpla el plan de trabajo.
\item[Dise�ador] Ingeniero de T�cnico de Telecomunicaciones con conocimientos en el �rea de telem�tica. Ser� el encargado de realizar el dise�o de los m�dulos del proyecto y los planes de pruebas de validaci�n.
\item[Programador] Programador con amplia experiencia de programaci�n en C y gran conocimiento del kernel Linux. Se encargar� de desarrollar los prototipos necesarios para realizar la validaci�n funcional del dise�o de los m�dulos.
\item[Administrativo] se encarga de resolver las cuestiones administrativas y burocr�ticas as� como de la documentaci�n del proyecto.
\end{description}

\subsection{Documentaci�n}
A continuaci�n se describen los documentos que deben utilizarse para la realizaci�n del proyecto. A pesar de ser un recurso material, se describe en un apartado diferente dada su importancia.

\begin{description}
\item [Documento de Memoria]
En este documento se exponen las razones que impulsan la realizaci�n de este proyecto. As�, en �l se recogen los objetivos que persigue este proyecto as� como los diferentes beneficios que reporta de ejecuci�n. Se enumeran las diferentes alternativas que se dan para la realizaci�n del mismo y se escogen las opciones que se consideran adecuadas para conseguir un dise�o �ptimo y un cumplimiento eficiente de los objetivos. 

Tambi�n se incluye un breve resumen y las especificaciones funcionales de los m�dulos del sistema. Para finalizar se adjunta un resumen del plan de trabajo y del presupuesto del presente proyecto.

\item [Documento de Dise�o de Alto nivel]
Este documento describen las funcionalidades y especificaciones de alto nivel de las diferentes partes del proyecto, sirve como gu�a general para saber cual es el razonamiento de dise�o que se ha seguido.

\item [Documento de Dise�o de Bajo nivel]
Este documento ofrece, desde un punto de vista t�cnico, una descripci�n detallada del modo en que se ha de implementar el proyecto. Se explica la forma que toma la soluci�n general explicada en el documento de Alto nivel, as� como ejemplos de c�digo que sirvan al Programador para guiarse.

\item [Documento de Pliego de condiciones]
El presente documento.
\end{description}


\chapter{Paquetes de trabajo}
A continuaci�n se describen los paquetes de trabajo a desarrollar en el presente proyecto.

\begin{description}
\item[PT1] \textbf{Estudio de las necesidades del sistema.}\\
Se realiza un estudio del entorno en el que se va a utilizar el sistema y de los usuarios potenciales del mismo. El objetivo es que las especificaciones queden perfectamente definidas.

\item[PT2] \textbf{Estudio de las alternativas.}\\
En este paquete de trabajo se identifican partes del sistema para las que se presentan diferentes alternativas de dise�o. Se realiza un estudio de dichas alternativas para cada caso y se escogen las opciones adecuadas. Es totalmente necesario para no reimplementar funcionalidades existentes.

\item[PT3] \textbf{Dise�o de los m�dulos del proyecto.}\\
Teniendo en cuenta las especificaciones planteadas y las alternativas escogidas, en este paquete de trabajo se realiza el dise�o de los diferentes m�dulos que componen el proyecto.

\item[PT4] \textbf{Desarrollo de los m�dulos del proyecto.}\\
En este paquete de trabajo se implementan una serie de m�dulos prototipo cuyo objetivo es validar el dise�o de los m�dulos dise�ados en el paquete de trabajo anterior.

\item[PT5] \textbf{Integraci�n de los m�dulos, pruebas de validaci�n y de rendimiento.}\\
Se ejecutan los prototipos de los m�dulos para su validaci�n y se realiza un estudio de los resultados obtenidos para validar el dise�o de los m�dulos.

\item[PT6] \textbf{Gesti�n del proyecto.}\\
Este paquete de trabajo contiene las tareas administrativas y de gesti�n que se deben realizar a lo largo de todo el proyecto as� como la generaci�n de la documentaci�n del mismo.
\end{description}

\chapter{Condiciones de recepci�n y pruebas}
n este apartado se definen las condiciones que se deben verificar en la entrega del proyecto. Se han definido una serie de pruebas que deben ser realizadas para proceder a la verificaci�n de los desarrollos y comprobar el correcto cumplimiento de las especificaciones recogidas en el documento de Memoria.

Las pruebas aqu� recogidas constituyen la base para la aceptaci�n del sistema, dejando al criterio del personal encargado del desarrollo la elaboraci�n de otras pruebas adicionales.

\section{Derechos de explotaci�n sobre el producto}
La herramienta que se desarrolle como resultado del proyecto ser� propiedad del dise�ador, Javier Domingo. La propiedad intelectual y el derecho al uso de la herramienta que se desarrolle, pertenecer� exclusivamente a Javier Domingo, respet�ndose el derecho moral del equipo investigador a ser designado como co-autores de los desarrollos.

El dise�o y el proyecto en general se ha licenciado con una licencia CC-BY-NC-SA. Se puede consultar esta licencia �ntegramente en \cite{CBSN13}. Las partes de c�digo que se han utilizado a modo explicativo y el c�digo asociado al mismo est� licenciado bajo la licencia GPLv3 \cite{GPLV3}.

El cliente se puede reservar el derecho a no distribuir el proyecto o el m�dulo, pero en caso de hacerlo, deber� hacerlo conforme a las citadas licencias.

\section{Entorno de pruebas}
Las pruebas se realizar�n en el marco del proyecto de continuaci�n planteado. La ubicaci�n de las pruebas, por tanto est� todav�a por definir.

\subsection{Materiales necesarios}
Dentro de los materiales necesarios para la realizaci�n de las pruebas se distinguir� el material hardware del software. En cuanto al hardware, los elementos necesarios son los siguientes:
\begin{itemize}
\item Un ordenador con una DAG \cite{DAPR10} que se utilizar� como inyector.
\item Un ordenador personal que se utilizar� para capturar las trazas inyectadas.
\item Un switch que conecta ambos ordenadores.
\end{itemize}

En lo relativo al software, ser�n necesarios dos sistemas Debian GNU/Linux, donde uno de ellos controlar� el inyector hardware que inserta las trazas de tr�fico configuradas, y el otro capturar� las trazas inyectadas para su posterior an�lisis.

\subsection{Personal necesario}
En lo referente al personal, ser�n necesarios los siguientes perfiles:
\begin{itemize}
\item Un director de proyecto que supervise el desarrollo de las pruebas.
\item Un ingeniero de dise�o que verifique el dise�o del sistema.
\item Un programador que comprueba el correcto funcionamiento del software.
\end{itemize}

\section{Pruebas de validaci�n}
En este apartado se especifican las pruebas a las que debe ser sometido el prototipo para la verificaci�n de su correcto funcionamiento y la validaci�n de los trabajos realizados. Como ya se ha indicado anteriormente estas pruebas se realizar�n sobre el prototipo que se implementa en el proyecto de continuaci�n propuesto. En dicho proyecto se ha incluido en el plan de trabajo un paquete de trabajo para realizar estas pruebas.


\subsection{Prueba del interfaz de usuario}
\subsubsection{Prop�sito}
Comprobar que el interfaz de los m�dulos hacia el usuario es capaz de mostrar datos por el sistema de ficheros procfs. El usuario deber� adem�s, poder cargar y descargar el m�dulo, y los datos deber�n ser reseteados en el proceso autom�ticamente.
\subsubsection{Procedimiento}
Se iniciar� la m�quina, y se cargar�n los m�dulos uno a uno comprobando 

Conclusi�n de la prueba:

La prueba se dar� por finalizada cuando se hayan mostrado todas las opciones de
ejecuci�n que presenta el sistema y se haya comprobado que se solicitan los par�metros
necesarios para cada caso y dichos par�metros son recogidos y almacenados
correctamente.
\subsection{Prueba del manager}
Prop�sito:
Comprobar que le manager realiza las llamadas adecuadas en funci�n de la informaci�n
que recibe. Y, por otra parte, comprobar que en cada llamada pasa correctamente los
par�metros adecuados.
Procedimiento:
Se inicia el manager y se le van pasando los c�digos de operaciones correspondientes y
los par�metros emulando la actuaci�n del interfaz de usuario. El manager por su parte,
ir� realizando llamadas pasando los par�metros adecuados.
Conclusi�n de la prueba
La prueba se da por concluida cuando se comprueba que el manager identifica
correctamente las operaciones que se pueden realizar con las llamadas que tiene que
hacer y los par�metros que tiene que pasar en cada caso.
\subsection{Pruebas del generador de datos}
Objetivo:
Comprobar que el generador de datos obtiene en una array los datos de inyecci�n seg�n
las distribuciones solicitadas para cada caso.
Procedimiento:

Introducir los par�metros indicados para la generaci�n de los datos y comprobar que los
datos salen correctamente almacenados en un array. Se deben comprobar todas las
distribuciones posibles.
Conclusi�n de la prueba:
La prueba concluye cuando se comprueba que el generador de datos es capaz de generar
datos de todas las distribuciones posibles de manera correcta. Almacenarlos en un array
de estructuras correctamente declarado y de tama�o indicado.

\subsection{Prueba del m�dulo de gesti�n}
Prop�sito:
Comprobar la correcta generaci�n de las trazas de inyecci�n: paquetes y datos de
inyecci�n.
Procedimiento:
En primer lugar, se le pasa al m�dulo de kernel los datos de inyecci�n desde espacio de
usuario. Se comprueba que han sido correctamente almacenados en espacio de kernel.
El m�dulo generar� los paquetes de datos. Hay que comprobar que se han generado de
manera correcta y que est�n correctamente almacenados. Tambi�n hay que comprobar
que el array de datos de gesti�n est� correctamente inicializado, actualizado y que es
accesible desde otros m�dulos de kernel.
Conclusi�n:
Esta prueba no finaliza hasta que no se generen flujos suficientes para agotar la
memoria de kernel y que el m�dulo de gesti�n tenga que avisar de que no hay espacio,
antes de machar informaci�n.

\subsection{Prueba de m�dulo de inyecci�n}
Prop�sito:

Comprobar que se realiza de manera correcta una inyecci�n.
Procedimiento:
Se programa una inyecci�n desde espacio de usuario pasando los par�metros adecuados.
Se comprueba que se realiza un acceso correcto a los arrays correspondientes para
obtener la informaci�n. Para esta prueba, hay que capturar la informaci�n inyectada en
otra m�quina para comprobar que el tr�fico est� correctamente modelado.
Conclusi�n de la prueba:
Esta prueba no termina hasta que no se compruebe el correcto funcionamiento del
inyector.

\subsection{Prueba de validaci�n global}
Prop�sito:
Comprobar el correcto funcionamiento del prototipo y la correcta interacci�n de todos
los m�dulos.
Procedimiento:
Se comprobar�n todas las funcionalidades del sistema. En este caso, cuando se decida
inyectar se capturar� en destino para comprobar la correcci�n de las distribuciones de
tr�fico y de tama�o de paquetes y la correcta temporizaci�n. Es importante comprobar
tambi�n la caracter�stica de inyecci�n por varios interfaces con varios inyectores
operando a la vez.
Conclusi�n:
La prueba concluye cuando se comprueba que el funcionamiento del prototipo es
estable.

\chapter{Condiciones econ�micas}
El presupuesto para la realizaci�n del proyecto asciende a un total de cincuenta y dos mil doscientos cuarenta euros (52.240?) impuestos no incluidos. Los trabajos adicionales a las labores establecidas en el presente contrato se comunicar�n al cliente y se facturar�n por separado, siempre que el proyectista lo considere oportuno.

\section{Forma de pago}
Se realizar� un primer pago por parte de la empresa contratante a la firma del contrato del presente proyecto, por valor de 6.000 ?, en concepto de conformidad con el mismo. Posteriormente, se realizar� un segundo pago por un importe total de 23.170 ?, en un plazo de cinco meses desde la recepci�n del primer pago. Finalmente, se realizar� un tercer pago final por un importe total de 23.170 ?, a la recepci�n de los trabajos realizados.

\chapter{Condiciones legales y contractuales}

\section{Actas de Recepci�n}
La recepci�n provisional de los servicios se efectuar� una vez prestados los mismos, debiendo firmar ambas partes el Acta de Recepci�n Provisional. A partir de la creaci�n de dicha acta, comenzar� el plazo de reclamaci�n, cuya duraci�n ser� de dos semanas. Transcurrido este plazo, y a no ser que el contratista notifique los defectos para su subsanaci�n dentro del plazo, se suscribir� el Acta de Recepci�n Definitiva.

\section{Responsabilidades del cliente}
El cliente se responsabiliza del cumplimiento de toda normativa legal vigente relativa a los servicios contratados, incluyendo los permisos y licencias de utilizaci�n, as� como la normativa de propiedad intelectual del desarrollo realizado.

\section{Responsabilidad del proyectista}
El presente proyecto requiere de una implementaci�n en el kernel actual para que se cumplan las especificaciones sobre las que se ha planeado el dise�o de bajo nivel. En caso de haber cambios importantes en versiones posteriores, habr�a que volver a comprobar que todas las especificaciones son suficientes.

\section{Extinci�n del contrato}
El contrato se extinguir� bien por conclusi�n o cumplimiento del mismo, o por resoluci�n de una de las partes involucradas en el mismo. Ser�n consideradas causas de resoluci�n:
\begin{itemize}
\item El incumplimiento de las cl�usulas contenidas en el pliego de condiciones.
\item La extinci�n de la personalidad jur�dica de la sociedad mercantil de una de las partes, salvo que el patrimonio sea incorporado a otra entidad.
\item Mutuo acuerdo entre las partes.
\item La declaraci�n de quiebra o suspensi�n de pagos de una de las partes.
\end{itemize}

\section{Resoluci�n de conflictos}
Los litigios que puedan surgir sobre interpretaci�n o modificaci�n del contrato ser�n resueltos por los Juzgados y Tribunales, renunciando a cualquier otro fuera que pudiera corresponder a las partes. Por tanto, se encomienda a estos el nombramiento de �rbitros y la administraci�n del arbitraje cuyo Laudo las partes se obligan a aceptar.
\chapter{Aspectos jur�dicos}
Se contemplan los siguientes aspectos jur�dicos relacionados con la contrataci�n y desarrollo del proyecto.

\section{Fuerza mayor}
El contratista no ser� considerado responsable por el incumplimiento de las obligaciones en tanto la ejecuci�n de los trabajos se retrase o se hiciese imposible por causas de fuerza mayor. Son considerados causas de fuerza mayor todos aquellos sucesos o circunstancias, fuera del control del contratista o del comprador y cualesquiera otras circunstancias que fueran imprevisibles, o que siendo previsibles fueran inevitables, de acuerdo con la jurisprudencia y la doctrina legal sentada sobre este concepto en el C�digo Civil.

Las antedichas causas de fuerza mayor se tomar�n en consideraci�n �nicamente cuando afecten directamente al suministro.

\section{Arbitrajes y tribunales}

Tanto el comprador como el contratista se comprometen a cumplir las condiciones establecidas en el presente Pliego de Condiciones y en la documentaci�n complementaria del contrato, resolviendo por medio de acuerdos y negociaciones las posibles diferencias que puedan surgir entre ellos respecto a la aplicaci�n, desarrollo, cumplimiento, ejecuci�n o interpretaci�n de los mismos.

En caso de que cualquier posible discrepancia o controversia entre ellos no pudiese ser llevada a buen fin, resolvi�ndose con �xito en la forma anteriormente indicada, el contratista se compromete a someter tales diferencias a arbitraje, formalizado de acuerdo a las normas reguladoras del mismo, contenidas en la vigente Ley de Arbitraje de Derecho Privado.

El arbitraje se realizar� en Bilbao y en la escritura ha de figurar el compromiso del plazo de 30 d�as como t�rmino en que los �rbitros han de pronunciar el Laudo correspondiente.
El procedimiento de arbitraje estar� de acuerdo con la Legislaci�n Espa�ola y ser�
resuelto en derecho por tres �rbitros, uno elegido por cada parte y el otro ya designado.
Si las partes no llegaran a un acuerdo para designar el primer �rbitro, �ste ser�
seleccionado por insaculaci�n entre los tres �rbitros de una lista que haya sido
proporcionada por el Colegio de Abogados de Bilbao.
Los �rbitros proceder�n a liquidar los gastos del procedimiento arbitral, incluyendo sus
honorarios, y a establecer a cu�l de las partes tendr�n que cargarse los gastos de
enjuiciamiento.
Las partes declaran desde ahora aceptar las decisiones del Colegio Arbitral.



\pagenumbering{Roman}

\bibliography{../bibliography}

\printindex

\end{document}
